% 等式约束 Lagrange 几何：等高线 + 约束曲线 + 梯度共线
\documentclass[tikz,border=8pt]{standalone}
\usepackage{ctex}
\usepackage{amsmath}
\usepackage{pgfplots}
\pgfplotsset{compat=1.18}
\usetikzlibrary{calc,arrows.meta,decorations.markings}

\begin{document}
\begin{tikzpicture}[scale=1.1]

%% 等高线族：f(x,y) = x^2/4 + y^2，椭圆
\foreach \r/\lab in {0.6/{}, 1.2/{$f{=}c_1$}, 1.8/{$f{=}c_2$}, 2.4/{$f{=}c_3$}} {
  \draw[gray!50, thin] (0,0) ellipse ({\r*2} and \r);
}
\node[gray, font=\scriptsize] at (2.55, 0.08) {$f{=}c_1$};
\node[gray, font=\scriptsize] at (3.65, 0.08) {$f{=}c_2$};
\node[gray, font=\scriptsize] at (4.85, 0.08) {$f{=}c_3$};

%% 约束曲线：圆 g(x,y) = x^2 + y^2 - 4 = 0（半径2）
\draw[blue!70!black, thick] (0,0) circle (2);
\node[blue!70!black, font=\small] at (1.6, 1.7) {$g(\mathbf{x})=0$};

%% 坐标轴
\draw[->, gray!60, thin] (-5.2,0) -- (5.5,0) node[right, font=\scriptsize] {$x$};
\draw[->, gray!60, thin] (0,-2.8) -- (0,2.8) node[above, font=\scriptsize] {$y$};

%% 最优点：等高线与圆相切于 (x*,y*) = (2*sqrt(2)/sqrt(5), 2/sqrt(5))
%% 取 f=x^2/4+y^2，约束 x^2+y^2=4
%% 梯度共线：(x/2, 2y) = lambda*(2x, 2y) => x/2/(2x) = 1/(4) = lambda => y/y=1, lambda=1/4
%% x*=2*cos(t), y*=2*sin(t), 代入: cos(t)/2*2*sin(t) = 1/(4)*(-2*sin(t),...
%% 简化：取最优点 (sqrt(2), sqrt(2)) 处验证（近似）
%% 精确计算：min x^2/4+y^2 s.t. x^2+y^2=4
%% Lagrange: x/2=2*lambda*x, 2y=2*lambda*y => lambda=1/4 (if x≠0), lambda=1 (if y≠0)
%% 不一致，取 x=0: y=±2; 取 y=0: x=±2
%% y=0,x=2: f=1; x=0,y=2: f=4 => 最小值在 (±2,0), 最大值在 (0,±2)
%% 切点在 (2,0): 等高线的切线是竖直方向，圆的切线也是竖直方向 => 相切 ✓

%% 最优点 P* = (2, 0)（极小）
\coordinate (Pstar) at (2, 0);
\fill[red!70!black] (Pstar) circle (2.5pt);
\node[below right, font=\small, red!70!black] at (Pstar) {$\mathbf{x}^*$};

%% 在 P* 处画梯度 nabla f = (x/2, 2y)|_{(2,0)} = (1, 0)（向右）
\draw[->, very thick, red!80!black, >=Stealth]
  (Pstar) -- ++(1.0, 0) node[right, font=\small] {$\nabla f$};

%% 在 P* 处画梯度 nabla g = (2x, 2y)|_{(2,0)} = (4, 0)，与 nabla f 共线
\draw[->, very thick, blue!70!black, >=Stealth, dashed]
  (Pstar) -- ++(0.6, 0) node[above, font=\small, blue!70!black] {$\nabla g$};

%% 标注梯度共线
\node[font=\small, align=center, red!60!black] at (3.2, -0.6)
  {$\nabla f = \lambda^* \nabla g$\\（梯度共线）};

%% 另一候选点 (-2, 0)（也是驻点，但是极大值）
\coordinate (Pmax) at (-2, 0);
\fill[black!50] (Pmax) circle (2pt);
\node[below left, font=\scriptsize, black!60] at (Pmax) {（极大）};

%% 一般点上的梯度（非切点处不共线的示意）
%% 在圆上 (1.41, 1.41) = (sqrt(2), sqrt(2)) 处
\coordinate (Q) at (1.41, 1.41);
\fill[black!40] (Q) circle (1.5pt);
%% nabla f at Q = (1.41/2, 2*1.41) = (0.71, 2.83)，方向
\draw[->, thick, orange!80!black, >=Stealth]
  (Q) -- ++(0.3, 1.0) node[right, font=\scriptsize, orange!80!black] {$\nabla f$};
%% nabla g at Q = (2*1.41, 2*1.41) = (2.83, 2.83)，方向 (1,1)
\draw[->, thick, cyan!60!black, >=Stealth]
  (Q) -- ++(0.55, 0.55) node[above, font=\scriptsize, cyan!60!black] {$\nabla g$};
\node[font=\scriptsize, orange!60!black, align=center] at (0.2, 2.4)
  {非最优点\\（梯度不平行）};

%% 约束曲线的切线方向（在最优点 P* 处，竖直线）
\draw[blue!50!black, thick, dashed] (2, -0.7) -- (2, 0.7);
\node[font=\scriptsize, blue!50!black, right] at (2.05, 0.55) {约束切线};

%% 标题
\node[font=\small\bfseries, align=center] at (0, -2.5)
  {等式约束极值几何：等高线族（椭圆）+ 约束曲线（圆）\\最优点处 $\nabla f \parallel \nabla g$（等高线与约束曲线相切）};

\end{tikzpicture}
\end{document}
